problem:  
Let \( (M^n,g) \) be a compact, simply connected Riemannian manifold with positive sectional curvature.  
For each unit vector \(v\in T_pM\), the **Jacobi operator** \(R_v\colon v^\perp\to v^\perp\), defined by \(R_v(w)=R(w,v)v\), has positive eigenvalues \(\lambda_1(v)\le\cdots\le\lambda_{n-1}(v)\).  
Define the **Jacobi spectral pinching constant**
\[
\Lambda(M)=\inf_{p\in M}\inf_{|v|=1}\frac{\lambda_1(v)}{\lambda_{n-1}(v)}\in(0,1].
\]
When \(\Lambda(M)=1\), \(M\) has constant sectional curvature.

It is known that if \(R_v\) is independent of \(v\) at each point (the **Osserman condition**), then under suitable hypotheses \(M\) must be locally two-point homogeneous.  
However, quantitative stability for near‑Osserman conditions is uncharted.

**Problem:**  
Does there exist, for each \(n\), a constant \(0<\eta_n<1\) such that if a positively curved, simply connected \(n\)-manifold satisfies  
\[
\Lambda(M)>\eta_n,
\]
then \(M\) is homeomorphic (or diffeomorphic) to a compact rank‑one symmetric space?  
Equivalently, can there exist a sequence of Riemannian manifolds \((M_i^n,g_i)\) with positive sectional curvature and \(\Lambda(M_i)\to1\), yet with \(M_i\) not homeomorphic to any CROSS?  

Why is it a "good" problem:  
This problem sits at the confluence of the **Osserman conjecture**, **pinching rigidity**, and **curvature operator geometry**.  Instead of the usual sectional curvature ratio, it measures geometric uniformity through the **spectral distribution of the Jacobi operator**, capturing how the geometry of geodesic deviations varies with direction.  The question asks whether near‑spectral‑symmetry (a quantitative almost‑Osserman condition) forces topological or smooth rigidity.  

A positive resolution would establish a powerful “spectral stability” bridge between curvature algebra and global topology, extending the Osserman‑rank‑rigidity philosophy into the realm of pinching.  A counterexample sequence would reveal new flexible regimes of almost‑homogeneous positive curvature.  
The definition is simple, but its implications reach deep into comparison geometry, curvature flow, and manifold classification—thus this is a truly *good* problem: conceptually minimal, topologically profound, and positioned just beyond today’s frontiers.